% Generado por gen_tabla_prev_modelo.py. No editar a mano.

\begin{table}[H]
    \centering
    \begingroup
    \scriptsize
    \renewcommand{\texttt}[1]{{\ttfamily #1}}
    \ttabbox[\FBwidth]{
        \caption{Prevalencia de \textit{Sí} (\%) de cada modelo por variable, en el nivel B1 ($N=1.313$). La columna \textbf{Prev. Sí} (sombreada) es la prevalencia anotada por las expertas, que sirve de referencia. En negrita, el modelo cuya prevalencia predicha más se acerca a la anotada, lo que no implica que marque las mismas piezas.}
        \label{tab:iris-prev-modelo-var}
    }{
        \begin{tabularx}{\textwidth}{
            >{\hsize=2\hsize\raggedright\arraybackslash\ttfamily}X
            >{\hsize=0.7875\hsize\columncolor{gray!12}\centering\arraybackslash}X
            >{\hsize=0.7875\hsize\centering\arraybackslash}X
            >{\hsize=0.7875\hsize\centering\arraybackslash}X
            >{\hsize=0.7875\hsize\centering\arraybackslash}X
            >{\hsize=0.7875\hsize\centering\arraybackslash}X
        }
            \toprule
            \normalfont\textbf{Variable} & \textbf{Prev. Sí} & \textbf{Gemini} & \textbf{GPT-4o-mini} & \textbf{GPT-5.4-nano} & \textbf{Gemma} \\
            \midrule
            lenguaje\_sexista & 81,0 & 3,9 & \textbf{46,2} & 0,1 & 1,2 \\
            masc\_generico & 81,0 & \textbf{66,9} & 19,1 & 9,7 & 36,9 \\
            sexismo\_discurso & 42,8 & 2,0 & 10,7 & \textbf{11,2} & 2,8 \\
            asimetria\_mujer\_hombre & 6,2 & 3,9 & 14,4 & 0,7 & \textbf{6,9} \\
            denominacion\_sexualizada & 10,0 & 0,9 & 1,2 & 0,1 & \textbf{1,4} \\
            \bottomrule
        \end{tabularx}
    }
    \endgroup
\end{table}
